% Results section for Kinetic AI paper.
% Numbers sourced from: research/memory/findings.md (F1–F25), results/exp01_mmd_vs_gda/ through results/scale/damage_sweep.json
% All statistics from at least 3 md5-distinct seeds except F25 (design measurement, single seed); every config hash and commit recorded.

\subsection{Magnetic Mirror Descent Convergence and Solver Engineering}
\label{sec:results:mmd}

\paragraph{Last-iterate convergence where simultaneous play cycles (F1).}

On symmetric zero-sum matrix games, fixed-magnet Magnetic Mirror Descent exhibits last-iterate linear (geometric) convergence while simultaneous gradient descent-ascent cycles indefinitely. The protocol: matching pennies and rock-paper-scissors, 2000 iterations per seed, 10 md5-distinct seeds, step size $\eta = 0.1$, regularization strength $\tau = 0.1$, uniform anchor policy (an equiprobable mixture). For each seed, the distance from the current iterate to the MMD fixed point was computed at every iteration, and a log-linear fit was applied to the final 50\% of the trajectory. MMD achieves $R^2 = 0.9948$ on matching pennies and $R^2 = 0.9015$ on rock-paper-scissors, both confirming linear decay. At the same step size and budget, simultaneous GDA (applied to the same payoff matrices, with separate gradients for row and column players) produces cycling dynamics with no monotone trend. GDA's final exploitability (NashConv) after 2000 steps is $1.93\ [1.90, 1.96]$ on matching pennies and $1.76\ [1.62, 1.88]$ on rock-paper-scissors, with $R^2 < 0.09$ across the final 50\% (no decay). Alternatives ruled out: initialization bias (all seeds start from uniform, identical across MMD and GDA), step size mismatch (fixed at $\eta = 0.1$ for both), and unequal budgets (both run 2000 steps). The linear convergence is specific to the magnetic anchor; removing the anchor term (setting $\tau = 0$) recovers GDA's cycling.

\paragraph{Asymmetric games and the context-dependent attractor (F2).}

For biased rock-paper-scissors (a 3×3 game with asymmetric payoffs), MMD-with-uniform-anchor converges to a deterministic attractor at NashConv $\approx 0.018$. However, the computed logit-QRE fixed point at $\lambda = 1/\tau = 10$ sits at NashConv $0.353$---the asymmetric attractor is distinct from the QRE. This observation is not explained by solver precision or averaging: the fixed-point identity ``MMD fixed point equals QRE($1/\tau$)'' holds exactly on symmetric games (where the uniform policy is already symmetric) but breaks on asymmetric games where the anchor is uniform but the equilibrium is not. The mechanism is context-dependent: the magnet and the game structure interact, producing an attractor that is neither NashConv $0$ (Nash equilibrium) nor the target logit-QRE. Practitioners transplanting MMD to asymmetric settings must verify convergence targets empirically.

\paragraph{Nash recovery via periodic reference resets (F3).}

Annealing the regularization strength---by re-centering the magnet on the current iterate at fixed intervals---recovers Nash equilibrium globally. The protocol: periodic reset of $y_{\text{ref}} \leftarrow y_t$ every 100 iterations on matching pennies, rock-paper-scissors, and biased rock-paper-scissors (same 10 seeds). Final NashConv: $8.48 \times 10^{-6}$ ($95\%$ CI $[4.78 \times 10^{-6}, 1.26 \times 10^{-5}]$) on matching pennies; $1.07 \times 10^{-6}$ on rock-paper-scissors; $5.21 \times 10^{-5}$ ($95\%$ CI $[2.62 \times 10^{-5}, 9.08 \times 10^{-5}]$) on biased rock-paper-scissors. The corresponding $R^2$ values range from $0.7181$ to $0.9825$ across the final 50\%, confirming smooth convergence. The reset protocol is the regularized Nash dynamics of Perolat et al., instantiated in a form reproducible on a single GPU. This rules out the alternative that the fixed-magnet result (F1) is merely a lower-dimensional accident; reset annealing generalizes to Nash.

\paragraph{QRE solver engineering (F7, F8).}

Two refinements accelerate and stabilize the QRE path-tracing algorithm. First, warm-starting: each fixed-point solve $s \leftarrow \mathrm{softmax}(\lambda A s)$ is initialized from the previous $\lambda$'s solution rather than uniform. On asymmetric $2 \times 2$ games, this reduces total iterations by $25.2\%$ (from $5990$ to $4481$ mean) with $\lambda \in \log\text{space}(0.01, 100, 50)$. The gain on biased rock-paper-scissors is smaller ($2.6\%$), and matching pennies (where all feasible $\lambda$ reach similar exploitability) is flat as expected. Exploitability along the path is not globally monotone for these games; path strategy movement is smooth and small (mean strategy shift $< 0.015$ in $\ell_2$ distance), ruling out the alternative that the path zigzags and the warm-start gain is illusory.

Second, plain fixed-point iteration $s \leftarrow \mathrm{softmax}(\lambda A s)$ diverges for $\lambda \approx 0.32$ onward on biased rock-paper-scissors. The iteration is a fixed-point map without contraction at high $\lambda$; the spectral radius at $\lambda = 0.32$ approaches 1. An adaptively damped variant $(1-\gamma)s + \gamma \cdot \mathrm{softmax}(\lambda A s)$ with adaptive step $\gamma = 1/(1 + \lambda/10)$ (initialized, then halved each time the residual increases) restores convergence across $\lambda \in \{1, 10, 100\}$ with iteration counts $21$, $700$, and $42{,}000$ respectively. The mechanism is classical (Picard iteration with damping); the discovery is that high-rationality QRE solves require damping by default at small-LM scale. The alternative (raw iteration) is now known to fail.

Figure~\ref{fig:mmd_convergence} visualizes matching pennies (the geometric-decay case for MMD) and the cycling of simultaneous GDA on the same problem.

\subsection{Mechanism Validation: Token Auction Truthfulness}
\label{sec:results:auction_mech}

\paragraph{Second-price token auctions are exactly truthful (F6).}

A second-price auction with per-token valuations set equal to each model's maximum next-token probability is empirically exactly truthful: the regret for truthful bidding is $0.0$ with 95\% CI $[0.0, 0.0]$. Protocol: 10 seeds, 200 auction trials per seed, agent configurations $\{3, 5\}$ agents, misreport grid $v \in \{0.25v, 0.5v, 1.0v, 1.5v, 2.0v\}$ (five multipliers of the true valuation), total observations $16{,}000$. For each trial, each agent bids one value on the grid; the mechanism computes the payment as the second-highest bid and scores the outcome. Regret is defined as the winner's payoff under the chosen bid minus its payoff under truthful bidding, summed over all agents. The empirical distribution of per-agent regret is concentrated at zero across all seeds and misreport levels. This confirms Vickrey's theorem in a practical LLM setting and validates the auction as the substrate for decoding experiments. The alternative (no mechanism check, assume Vickrey applies by theory alone) is ruled out by direct measurement.

A weighted-logit-average aggregation baseline (mixing model outputs as $p \propto w_i \sigma(v_i)$ where $w_i$ are learned weights) is measurably manipulable: mean misreport gain is $0.0773$ at $n = 3$ and $0.0683$ at $n = 5$. Its payments are not VCG. This establishes second-price auctions as the appropriate choice for truthful aggregation in multi-model decoding.

\subsection{EqLM at Smoke Scale and First Diagnostics}
\label{sec:results:eqlm_smoke}

An end-to-end pipeline was validated at smoke scale (800 steps, 1000 BabyLM pairs, 300 step evaluation budget). An initialization bug in the first run set all models to loss $\approx 125$ instead of $\ln|V| \approx 10.8$; correction revealed exact parameter-matched parity between the weight-tied EqLM and an explicit transformer: final loss $5.94$ (EqLM) vs.\ $5.99$ (explicit), with EqLM at $6.2\%$ lower parameter count ($10.67$M vs.\ $11.06$M). BLiMP scores at this stage ($0.452$--$0.465$ both arms) register as chance performance at small scale, consistent with a $3$-minute pretraining epoch. Two critical optimizer defects emerged during root-cause analysis. The magnetic anchor's strength at $\tau \leq 10^{-2}$ is loss-neutral in pretraining: final loss across magnetic arms ranged $8.52$--$8.61$ versus explicit $8.51$, within the pre-registered $10\%$ tolerance. Coupled weight decay in AdamW silently destroys sparse-gradient parameters by acting as an $\text{lr} \cdot \text{wd}$-magnitude decay on rows whose gradients are mostly zero: EqLM+MagneticAdamW at $\tau = 0$ reached only $10.46$ against the decoupled formulation's exact $8.8236$ on tied embeddings. These findings led to a methodological decision: pretraining uses plain AdamW; magnetic anchoring is reserved for preference-phase fine-tuning, its theoretical home.

\subsection{Convergence Crisis and the Fixed-Point Map Fix}
\label{sec:results:diagnosis}

The v1 residual-form map $f(z,x) = (1-\alpha)z + \alpha(z + \mathrm{Attn}(z) + \mathrm{MLP}(z) + x)$ admits no bona fide fixed point. Solver residuals plateau at constant magnitude with tail ratio $\approx 0.99$ over $100$ iterations; residual magnitude scales \emph{linearly} with damping coefficient $\alpha$ (residuals: $32.65 \to 5.41 \to 1.35$ at $\alpha = 1 / 0.2 / 0.05$), the signature of $z \leftarrow z + \alpha g(z)$ with $\lVert g \rVert$ constant. The v1 models were weight-tied iterative transformers, not equilibrium models; this reframing explains their performance.

The v3 post-LN map restores bounded iterates and genuine fixed-point convergence. Following DEQ-transformer practice, the map applies layer normalization inside the recursion: $h = \mathrm{LN}_1(z + \mathrm{Attn}(z) + x)$ and $f(z, x) = \mathrm{LN}_2(h + \mathrm{MLP}(h))$. Convergence is gated on the relative residual $\lVert f(z^*) - z^* \rVert / (\lVert z^* \rVert + \epsilon)$. At smoke scale with $d = 204$, $T = 127$, and $80$ iterations, this map exhibits bounded approach to a fixed point: relative residual decays from $1.0$ to $0.105$ monotonically, then creeps slowly ($0.1078 \to 0.1053$ over the last $24$ iterations). The fixed-drift signature of v1 is eliminated, yet the effective contraction factor remains $\approx 1$ at LM width, rendering tolerance $10^{-3}$ unreachable within practical iteration budgets.

A striking discovery makes the gap addressable. Adding an auxiliary loss term $\lambda_c \cdot \lVert f(z^*, x) - z^* \rVert / \lVert z^* \rVert$ applied only along the solver's trajectory (one block application after the no-gradient solve) reduces the exit relative residual by a factor of $16\times$ (from $0.1277$ to $0.0076$--$0.0078$) for every $\lambda_c \in \{0.01, 0.1, 1.0\}$ at a cost of $\leq 1.8\%$ in cross-entropy loss ($8.92$--$8.99$ vs.\ control $8.77$) and zero wall-time overhead ($0.99$--$1.00\times$ wall-clock). The model learns to become contractive during training; the equilibrium computation is taught, not built in by architecture alone.

\subsection{H1 Iteration 1 and 2: The 11M-Parameter Verdict}

The first full run (20k steps, BabyLM strict-small, 1000-pair BLiMP evaluation) missed the 95\% parity gate decisively. Explicit 12-layer baseline: BLiMP $0.734$. EqLM v1 (residual map, no contraction fix): $0.571$. EqLM v3 + post-LN (map fix alone): $0.584$. The v1 and v3 arms are not statistically separable ($\approx 0.8\sigma$ at $n = 1000$ pairs), and both fall far below the gate at $78$--$80\%$ of baseline.

With the post-LN map and solver-aware loss implemented, the second iteration improved substantially. On the first seed, EqLM reached $0.704$ (94.4\% of baseline's $0.746$), statistically a tie with the 95\% threshold. The mandatory 3-seed verdict across seeds 42, 43, and 44 landed at: A1 (explicit) $0.746 / 0.689 / 0.705$ (mean $0.7133 \pm 0.0294$); A3 (EqLM post-LN) $0.704 / 0.654 / 0.633$ (mean $0.6637 \pm 0.0365$). Per-seed ratio: $0.944 / 0.949 / 0.898$; mean ratio $0.9303$; bootstrap $95\%$ CI $[0.8979, 0.9492]$. The upper confidence bound falls below $0.95$, formally missing the pre-registered threshold. The progression from iteration 1 (ratio $0.78$) to iteration 3 (ratio $0.93$) documents the value of diagnosed fixes: initialization scaling and the post-LN map.

Warm-started equilibrium decoding cuts solver iterations dramatically. Initializing each decode step from the previous token's equilibrium cache reduces iterations per token from $7.91$ to $1.64$ (a $79.3\%$ reduction, exceeding the pre-registered $\geq 50\%$ gate). Greedy-token agreement reaches $97.6\%$ (976/1000 tokens), with 39 out of 40 complete sequences matching exactly, narrowly missing the pre-registered $\geq 99\%$ agreement gate. An explicit 12-layer stack always incurs 12 block applications per token; a warm-started equilibrium model requires only $\approx 1.6$ at this scale.

\subsection{The 121M-Parameter Study and Mechanism Correction}
\label{sec:results:scale}

At 121M parameters (batch 32, 10k steps, 3 seeds on RTX 5090), the gap reappeared under implicit-differentiation training. Explicit baseline: BLiMP $0.6833$ (seeds 42/43/44: $0.682 / 0.675 / 0.693$). EqLM v3: BLiMP $0.5377$ (seeds 42/43/44: $0.537 / 0.532 / 0.544$). Per-seed ratio: $0.7874 / 0.7881 / 0.7850$; mean $0.7868$; bootstrap $95\%$ CI $[0.785, 0.788]$; paired $t = -82.6$.

The original mechanistic claim---``the solver converges at small width and fails at large''---was refuted by adversarial review: solver telemetry shows convergence rate $0.0$ and all 12 iterations used at \emph{both} scales (11M and 121M). The corrected mechanism is graded, not binary: under a fixed 12-iteration budget, as width scales from $204$ to $1704$ dimensions, the fixed-point map becomes progressively less contractive per iteration, and the truncation penalty widens. At 11M, the quality gap is $\approx 7\%$; at 121M, $21\%$. This is a graded contraction-versus-width effect, not a binary solver failure.

Two cost facts survive audit. EqLM maintains a peak-memory advantage: $6.29$ GB versus $8.13$ GB for the explicit baseline ($-23\%$), consistent with $O(1)$ depth-memory scaling. Wall-clock cost is substantial: $92$ min/arm for EqLM versus $11$ min/arm for the explicit baseline ($8.4\times$ slower per matched step). The comparison confounds two differences: step budget (10k vs.\ 20k from the 11M run, applied identically to both arms) and width. The trend's direction is robust; its magnitude is confounded.

\subsection{Anytime-Unrolled Training Reaches Parity}
\label{sec:results:parity}

When the training regime changed from implicit-differentiation to anytime-unrolled supervision, the scale gap dissolved entirely. The same weight-tied block is unrolled explicitly for 12 iterations from $z_0 = x$, then supervised with cross-entropy loss at intermediate iterates $z_4, z_8, z_{12}$ (weights $0.15 / 0.3 / 1.0$). Evaluation still runs the standard Anderson solver at budget 12, so the trained and evaluated computation paths differ by algorithm but operate at the same iteration cap. Results across seeds 42, 43, 44: BLiMP $0.662 / 0.697 / 0.672$ (mean $0.677$, ratio $0.991$). Seed 43 exceeds its baseline ($0.697$ vs.\ $0.675$).

Recipe audit confirms the comparison is clean: identical token-cache file (same data stream), steps ($10$k), batch ($32$), learning rate ($3 \times 10^{-4}$), weight decay ($0.01$), gradient clipping ($1.0$), and parameter count within $2.5\%$ (120.7M EqLM vs.\ 123.8M explicit). Three dividends emerge. First, wall-clock is reversed: unrolled training is $2.1\times$ \emph{faster} than implicit-differentiation training ($44$ vs.\ $92$ min/arm). The fixed-point solve, not backpropagation, dominates EqLM's computational cost. Second, the pre-registered budget-sweep rider is met: the same checkpoint evaluated at Anderson budgets $\{4, 8, 12\}$ yields BLiMP scores that degrade gracefully. Seed 42: $0.621 / 0.638 / 0.662$. Seed 43: $0.601 / 0.674 / 0.697$. Seed 44: $0.587 / 0.655 / 0.672$. One model serves three compute budgets; the adaptive-computation property is realized through the architecture itself. Third, a real cost: unrolled training stores its iterates for backpropagation, incurring activation memory equivalent to the explicit model ($16.4$ GB vs.\ $7.8$ GB during training). The serving-time properties---$O(1)$ depth-memory, warm-start decoding---are unchanged because evaluation still runs the implicit solver.

\begin{table}[t]
\centering
\small
\caption{Pretraining results at 121M parameters on BabyLM strict-small (batch $32 \times 10$k steps, 3 seeds, 1000-pair BLiMP subset). Ratio is per-seed EqLM / explicit. The anytime-unrolled training regime recovers the gap solver-based training had worsened at width.}
\label{tab:main_results}
\begin{tabular}{@{}lccccc@{}}
\toprule
Model (regime) & BLiMP (per seed) & Mean & Ratio & Peak mem & Train time \\
\midrule
Explicit 12-layer (baseline) & 0.682 / 0.675 / 0.693 & 0.683 & --- & 8.1 GB & 11 min \\
EqLM (implicit, F20) & 0.537 / 0.532 / 0.544 & 0.538 & 0.787 & 6.3 GB & 92 min \\
EqLM (anytime-unrolled, F24 B1) & 0.662 / 0.697 / 0.672 & 0.677 & \textbf{0.991} & 16.4 GB$^{\dagger}$ & 44 min \\
\bottomrule
\multicolumn{6}{@{}l@{}}{\footnotesize $^{\dagger}$Training-phase cost of storing unrolled iterates; serving uses the implicit solver and 6.3 GB.}
\end{tabular}
\end{table}

\subsection{Certification and the Quality-Preservation Challenge}
\label{sec:results:certification}

H6 demanded both quality recovery (closing $\geq 50\%$ of the width gap) and certification (solver convergence rate $> 0.5$ at evaluation) in a single model. A TRIZ-style problem analysis framed this as a physical contradiction: the fixed-point map must be contractive (small Lipschitz constant, for certification) and expressive (large capacity, for quality). Spatial separation---penalizing contraction \emph{only along the solver trajectory} rather than globally---was proposed as a resolution.

Arm B2 implements trajectory-local contraction via a finite-difference probe of the map at a uniform-random point on the solve ray, one probe per training step. The penalty term $\lambda_c \cdot \lVert f(z + \epsilon v) - f(z) \rVert / \epsilon$ (with $\epsilon = 10^{-3}$) is applied at evaluated iteration cost $\approx 0\%$ wall-clock. Result (seed 42, screen level): BLiMP $0.577$, above control but below parity. Solver convergence rate reaches the full value of $1.0$ at mean $4.0$ iterations, the program's first certified 121M-parameter equilibrium. Peak memory is lowest of all arms ($5.5$ GB).

Arm B3 applies spectral normalization to a shallow 256-dimensional core sandwiched between wide explicit encoder/decoder layers ($2 + 2$ layers at $d = 1152$). Result: BLiMP $0.642$, partial certification. Arm B4 combines anytime supervision (arm B1) with the raw trajectory penalty from B2, pre-registered with the prediction that this combination would meet both gates. The combination was refuted: the unnormalized penalty reached $3.5 \times 10^4$ at initialization and drowned the cross-entropy supervision signal under gradient clipping. Final BLiMP $0.529$ (below control), convergence rate $0.0$, final loss $\hat{L} \approx 2.5$--$3.9$. A documented failure of the naive sum; the lesson is sharp: quality-preserving certification remains the open problem. H6 closes as \textsc{partial}---its two halves (quality at parity, certification achieved) are separately solved, but their straightforward combination fails.

\subsection{Preference Optimization: Where Magnetic Anchoring Does Not Help}

H3 tested whether a magnetic proximal pull in parameter space improves preference optimization. The method, parameter-space magnetic anchoring (PMA), applies a magnetic term inside AdamW toward frozen base weights under the standard DPO loss---distinct from Wang et al.'s policy-space Magnetic Preference Optimization (MPO), which targets policy-space Nash equilibrium. The testbed makes the preference signal exact: BLiMP minimal-pair judgments ($\text{good} \succ \text{bad}$) serve as preference pairs with a phenomenon-level train/test split (600 training, 400 held-out unseen phenomena). Each seed fine-tunes its own 121M checkpoints from exp10; arms differ only in magnetic strength $\tau$.

Verdict: \textsc{partial}. At the pre-registered $\tau \in \{10^{-3}, 10^{-2}\}$, PMA held-out accuracy equals DPO \emph{exactly} (all seeds, both bases)---the $\geq$ DPO half holds by equality. The promised KL reduction does not appear: PMA versus DPO KL drift on the explicit base is $-0.0018 / -0.0009 / +0.0004$ (mean $-0.0008$, not significant). Total magnetic displacement scales as $\eta \tau T \approx 2 \times 10^{-5}$ relative to the base---under-dosed. A dose-response rider swept $\tau \in \{0.1, 1, 10\}$ with the pre-registered prediction that drift reduction becomes visible by $\tau = 1$. The prediction was refuted: even $\tau = 10$ moves KL only $1.2406 \to 1.2256$ ($-1.2\%$) with accuracy unchanged. The magnetic pull is second-order to the DPO gradient across four orders of magnitude, consistent with the pretraining finding that the magnet's natural function is stability during pretraining, not drift control during fine-tuning.

Two secondary findings survive. On the explicit base, standard DPO lifts trained-phenomena accuracy $0.646 \to 0.877$ while held-out (unseen) phenomena drop $0.740 \to 0.612$, with $1.24$ nats/token KL from base. Preference optimization damages out-of-distribution phenomena. The equilibrium architecture is intrinsically drift-resistant: under identical loss, learning rate, and training steps, EqLM moves train accuracy only $0.493 \to 0.516$ with $1655\times$ less KL drift ($0.00075$ nats/token vs.\ $1.24$). Audit ruled out the trivial explanation (frozen base already satisfied): EqLM's base train accuracy $0.493$ is low relative to the explicit base $0.646$. The cause is damped gradients through the 12-iteration weight-tied solve. This dual nature is reported as-is: stability against preference-induced drift and insensitivity to preference optimization at matched learning rate, two sides of the same architectural property.

\subsection{Token Auctions at Scoring Time}

H4 tested whether per-token auction selection over domain specialists beats any fixed model choice. Two 30M-parameter specialists per seed trained on disjoint BabyLM subdomains (child-directed speech, simple-Wikipedia text; line-level 5\% held-out split per seed). Evaluation on a $50/50$ interleaved held-out mixed stream ($\approx 100$k tokens). At each position, every specialist bids its own maximum next-token probability (target-independent by construction); the highest bidder's distribution scores the true token; the winner pays the second-price (verified earlier to be exactly truthful). Results (3 seeds, 42/43/44):

Mixed-domain perplexity. Auction: $158.5 / 189.4 / 200.5$ (mean $182.8$). Best single specialist: $234.3 / 232.8 / 243.4$ (mean $236.8$). Uniform logit-average ensemble: $207.9$ (interpolated mean). Worst specialist: $1242.4$ (interpolated mean). The auction beats the best single on every seed (paired $t = 4.98$). The decomposition clarifies the mechanism: confidence-bid selection is imperfect within a domain (child-speech windows where the auction scores $26$--$33$ ppl trail the resident specialist at $13.8$), yet per-token selection dominates any fixed commitment across mixed domains. Each specialist collapses off-domain to $\approx 4000$--$4900$ ppl; the auction avoids these cliff edges at every transition.

\subsection{Auction Decoding in Closed Loop: A Negative Result and Metric Lesson}
\label{sec:results:closedloop}

H5 tested whether the auction's advantage persists when each system generates its own context. Protocol: 100 mixed prefixes (50/50 childes/simple-wiki), 32 tokens generated per prefix, greedy decoding. Judge: frozen 124M explicit LM baseline, NLL/token on generated text only.

Verdict: \textsc{missed}. Judge NLL: best single specialist $3.39 / 3.37 / 3.69$ (per seed) versus auction $4.74 / 4.23 / 4.60$. The pre-registered advantage does not appear. Teacher forcing re-anchors per-token selection to the true context at every position; closed-loop generation removes this anchor. The auction drifts toward one specialist's style regardless of prompt domain, measured independently via domain consistency ($\approx 1.0$ on childes-prefixed continuations but $\approx 0.02$ on wiki-prefixed). The scoring-time scoping imposed by adversarial review before F22 ran is empirically vindicated here.

Deeper investigation reveals a metric pathology. A per-domain decomposition (reproduced bit-identically on two machines) shows the child-speech specialist scoring best \emph{even on wikipedia-only prompts} (judge NLL $3.27$--$3.78$), while the wikipedia specialist generating on its own domain scores $5.10$--$5.49$. The judge, trained on a child-speech-heavy BabyLM diet, carries a $\approx 2$-nat style prior that dwarfs every between-system effect. All systems were evaluated on the same stream; the metric is judge-relative, not model-invariant. Two observations survive independently of the judge: the auction's output measurably diverges from teacher-forced behavior, and the auction exhibits the lowest degeneration (3-gram repetition $0.387$--$0.483$) while uniform logit averaging collapses ($0.780$--$0.832$)---bid competition acts as an implicit anti-repetition regularizer. Exposure bias compounding across models is a documented follow-up.

\subsection{Conversion at 1.7B: The Damage Curve}

Conversion of Qwen3-1.7B (28 layers, $d = 2048$, 1.721B parameters) to the EqLM topology was measured without uptraining to inform the architecture design. The topology divides layers into explicit outer blocks (pre/post) and an implicit core (one to four iterating blocks). A sweep over outer-layer count ($n_{\text{pre}} = n_{\text{post}} \in \{6, 8\}$), core count ($n_{\text{cores}} \in \{1, 2, 4\}$), and initialization (average or stepwise layer averaging) revealed three findings.

Average-layer initialization preserves far more function than stepwise adoption of a single representative layer. For example, at 8+8 outer layers with 1 core: average initialization yields $1.9 \times 10^3$ ppl versus stepwise $2.2 \times 10^5$ ppl (two orders of magnitude difference). Collapsing a group to their mean preserves representational capacity; picking one representative destroys it.

Explicit outer layers matter more than the number of cores. Going from 6+6 to 8+8 outer layers (adding $8$ parameters to explicit layers) with 1 core improves ppl from $18{,}465$ to $1{,}909$ (10-fold). Going from 1 to 2 cores at fixed 8+8 outer layers changes ppl only $1{,}909$ to $1{,}933$ (negligible). The first and last layers carry function destroyed by tying; the middle layers are redundant.

Every configuration starts heavily damaged pre-uptraining ($300$--$3{,}000\times$ the base ppl of $6.01$). Conversion is not a free lunch; the uptraining budget, not the surgery, is the binding constraint.

The operating point was selected pre-uptraining: $n_{\text{pre}} = n_{\text{post}} = 8$, $n_{\text{cores}} = 1$, average initialization, yielding 1.167B parameters ($68\%$ of base) and starting ppl $1{,}909$. The pre-registered parameter gate was amended from $\leq 60\%$ to $\leq 70\%$ on evidence that the $56$--$59\%$ configurations start $3$--$10\times$ more damaged. Uptraining runs are in progress; results are not yet available.

Figure~\ref{fig:parity_budget} displays the anytime checkpoint's graceful degradation across inference budgets and the ratio of EqLM to baseline across the 11M and 121M studies.
